\subsection{Admissibility from Bounded Flux}
  \label{subsec:admissibility-bounded-flux}

  The dynamics of the $\chi$ substrate is constrained by a fundamental bound on
  the local rate of relaxation.
  This bound expresses the finite transmissibility of the relational substrate
  and limits the rate at which local configurations may evolve.

  Operationally, this constraint can be expressed as
  \begin{equation}
    |\partial_t \chi_v| \le c_\chi ,
  \end{equation}
  where $c_\chi$ denotes the maximal relaxation speed associated with the
  substrate.

  This condition plays a role analogous to a causal bound, but it operates at the
  level of relational relaxation rather than geometric propagation.
  Configurations whose local relaxation flux exceeds this bound cannot be
  sustained dynamically.

  When the dynamics is expressed in a spectral basis of the relational Laplacian
  $L$, each mode evolves independently at leading order.
  Let $\lambda_\rho$ denote the Laplacian eigenvalue associated with a spectral
  sector labelled by $\rho$.

  The bounded-flux condition then restricts the maximal amplitude that a spectral
  mode can reach without violating the local relaxation constraint.
  A detailed analysis shows that the maximal admissible amplitude of a mode
  associated with eigenvalue $\lambda_\rho$ obeys
  \begin{equation}
    A_\rho^{\max} = \frac{c_\chi}{\sqrt{\lambda_\rho}} .
  \end{equation}

  This relation defines a universal admissibility envelope for spectral modes.
  Modes associated with smaller eigenvalues can sustain larger amplitudes,
  whereas high-frequency modes are strongly suppressed by the bounded-flux
  constraint.

  In this sense, the Laplacian spectrum acts as a structural filter on the space
  of dynamically admissible configurations.
  Only modes compatible with the admissibility envelope remain dynamically
  sustainable.

  The existence of this envelope introduces a hierarchy among spectral sectors.
  Even when several sectors are kinematically available, the bounded-flux
  constraint favours those associated with lower Laplacian eigenvalues.

  The consequences of this spectral filtering become particularly transparent
  when explicit relational graphs are considered.
  The simplest analytically tractable example is provided by the Cayley graph of
  the quaternion group $Q_8$, which we analyse in the following section.
